% !TeX root = ../sections/02_exercises.tex
\begin{exercise}
	数列${a_n}_{n_\in\mathbb{N}}$と${b_n}_{n_\in\mathbb{N}}$が収束して$\lim_{n\rightarrow\infty}a_n=\alpha,\lim_{n\rightarrow\infty}b_n=\beta$であり、「ある番号$N\in\mathbb{N}$が存在して、$n\geq N$ならば$a_n\leq b_n$」が成立するとき、$\alpha\leq\beta$が成立することを証明せよ。
\end{exercise}
\begin{background}
	この問題では，十分大きな番号以降で成り立つ不等式が，
	極限にも受け継がれることを示す。
	
	\begin{definition}[数列の収束]
		数列 $\{a_n\}$ が実数 $\alpha$ に収束するとは，
		任意の $\epsilon>0$ に対して，ある自然数 $N$ が存在して，
		すべての $n\geq N$ に対して
		\[
		|a_n-\alpha|<\epsilon
		\]
		が成り立つことをいう。
		このとき，
		\[
		\lim_{n\to\infty}a_n=\alpha
		\]
		と書く。
	\end{definition}
	
	\begin{remark}
		今回の仮定は
		\[
		\lim_{n\to\infty}a_n=\alpha,
		\qquad
		\lim_{n\to\infty}b_n=\beta
		\]
		である。したがって，$n$ を十分大きくすれば，
		$a_n$ は $\alpha$ に近く，$b_n$ は $\beta$ に近くなる。
	\end{remark}
	
	\begin{remark}
		また，問題文では，ある番号 $N_0\in\mathbb{N}$ が存在して，
		すべての $n\geq N_0$ に対して
		\[
		a_n\leq b_n
		\]
		が成り立つと仮定されている。
		
		これは，すべての $n$ で $a_n\leq b_n$ が成り立つという意味ではなく，
		ある番号以降では常に $a_n\leq b_n$ が成り立つ，という意味である。
	\end{remark}
	
	\subsection{極限における順序保存性の準備}
	
	\begin{lemma}[有限個の自然数の最大値]
		自然数 $N_0,N_1,N_2$ に対して，
		\[
		N=\max\{N_0,N_1,N_2\}
		\]
		とおけば，
		\[
		n\geq N
		\]
		ならば
		\[
		n\geq N_0,\qquad n\geq N_1,\qquad n\geq N_2
		\]
		がすべて成り立つ。
	\end{lemma}
	
	\begin{remark}
		この補題により，
		$a_n\leq b_n$ が成り立つ条件と，
		$a_n$ が $\alpha$ に近い条件，
		$b_n$ が $\beta$ に近い条件を，同時に満たすような
		十分大きな $n$ を取ることができる。
	\end{remark}
\end{background}

\begin{strategy}
	示したいことは
	\[
	\alpha\leq \beta
	\]
	である。
	
	このような不等式を示すときは，背理法を用いると考えやすい。
	そこで，反対に
	\[
	\alpha>\beta
	\]
	であると仮定する。
	
	このとき，
	\[
	\alpha-\beta>0
	\]
	である。そこで，この正の差を利用して
	\[
	\epsilon=\frac{\alpha-\beta}{3}
	\]
	とおく。
	
	$a_n\to\alpha$ であるから，十分大きな $n$ に対して
	\[
	|a_n-\alpha|<\epsilon
	\]
	が成り立つ。これは特に
	\[
	\alpha-\epsilon<a_n
	\]
	を意味する。
	
	また，$b_n\to\beta$ であるから，十分大きな $n$ に対して
	\[
	|b_n-\beta|<\epsilon
	\]
	が成り立つ。これは特に
	\[
	b_n<\beta+\epsilon
	\]
	を意味する。
	
	ここで，
	\[
	\epsilon=\frac{\alpha-\beta}{3}
	\]
	と選んでいるので，
	\[
	\alpha-\epsilon
	>
	\beta+\epsilon
	\]
	が成り立つ。実際，
	\[
	\alpha-\epsilon-(\beta+\epsilon)
	=
	\alpha-\beta-2\epsilon
	=
	\alpha-\beta-\frac{2}{3}(\alpha-\beta)
	=
	\frac{1}{3}(\alpha-\beta)
	>0
	\]
	である。
	
	したがって，十分大きな $n$ に対して
	\[
	a_n>\alpha-\epsilon
	>
	\beta+\epsilon
	>b_n
	\]
	となる。つまり
	\[
	a_n>b_n
	\]
	が得られる。
	
	しかし，問題文の仮定より，ある番号 $N_0$ 以降では常に
	\[
	a_n\leq b_n
	\]
	が成り立つ。
	
	十分大きな $n$ を取れば，
	$a_n\leq b_n$ と $a_n>b_n$ が同時に成り立つことになり，
	これは矛盾である。
	
	よって，仮定
	\[
	\alpha>\beta
	\]
	は誤りである。したがって
	\[
	\alpha\leq\beta
	\]
	が成り立つ。
\end{strategy}